\section{Executive Summary}\label{sec:summary}

NTNU wants to explore how Generative AI (GenAI) and Large Language Models (LLMs) can improve its teaching processes. Today, these processes are supported unevenly by Studentweb, Canvas, Panopto and Inspera. Activities that store and distribute information, such as course registration, publishing recordings and conducting exams, are well supported. Activities that require interpretation and judgment remain largely manual.

This report models the AS-IS teaching process in BPMN and identifies three sub-processes where this gap is clearest. In Case 1, lecture recordings are stored as raw video, so every student must scrub through the recordings and write their own notes. In Case 2, feedback on exercises is often limited to a score, since whether a comment is written depends on TA workload. In Case 3, course selection offers no personalization, so students research electives manually and depend on study advisors with unknown response times.

The TO-BE solution adds one shared GenAI service to the existing systems. The service generates a transcript, a topic index and study material once per lecture. It drafts a score and detailed feedback for every submission. It also recommends courses after a rule-based eligibility check against the student record in FS. A human remains responsible for every output: lecturers approve the study material, TAs approve scores and feedback, and students and study advisors make the final course choices.

Reaching the TO-BE situation is mainly about changing the work around existing systems, not about building new ones. Case 1 is the least disruptive change, while Case 3 is the most demanding, since it requires new data flows from FS and up-to-date course descriptions. We recommend introducing each change through pilots structured with Lewin's change model, with the new functionality built into tools the staff already use. The SusAF analysis shows benefits in personal support, consistency and staff time. It also shows risks: personal data processed by an external service, reviewers approving AI output without reading it, overreliance on AI, dependency on one vendor, and higher energy use.

\section{Background}\label{sec:background}

\subsection{ICT in academic education}\label{sec:ict}

Digital technology has changed higher education in stages. Learning management systems (LMS) such as Canvas have centralized course material, communication and submissions. Lecture recording has added a second stage. Recordings of face-to-face lectures are made available online, and students make significant use of them, with some demonstrated benefits for learning outcomes, while lecturers worry about effects on attendance and engagement \parencite{ocallaghan2017}. Recordings also enable the flipped classroom, where students study material before class and apply it during class. A meta-analysis of 114 studies found a small positive effect on learning outcomes but no effect on student satisfaction \parencite{vanalten2019}.

Assessment has also been digitized. In Norway, LMSs such as Canvas and Blackboard are used for ungraded assessment tasks, while dedicated e-exam systems like Inspera Assessment are used for graded ones \parencite{chirumamilla2021}. At NTNU, digital exams grew from 282 submissions in 2016 to 68\% of all examinations in 2024 \parencite{inspera2025}.

Study administration has been digitized as well. At NTNU, course and exam registration is handled in Studentweb, the student-facing interface to the Common Student System (FS), which stores study programmes, completed courses and grades. Information about course content is published separately on the NTNU course pages and Innsida. These systems validate and store students' choices, but they do not help students decide what to choose; guidance on electives is given manually by study advisors.

Core parts of the teaching process are therefore supported by information systems, but the parts that require judgment, such as making lecture content usable, giving feedback and guiding course choices, remain largely manual. This work is either done by each student individually, as when every student reviews the same recording, or limited by staff capacity, as when TAs only write comments when there is time. These are the parts of the process this report focuses on, and they are analysed as three cases in Section~\ref{sec:asis-cases}.

\subsection{The Generative AI shift}\label{sec:genai}

Earlier systems store and distribute content. Generative AI (GenAI) and large language models (LLMs) can also interpret and produce it, for example by summarizing material, generating feedback and supporting learners. Reviews identify automated feedback and assessment, learning support and the development of critical skills as the main pedagogical uses. They also identify challenges: unreliable or biased output, privacy and academic integrity risks, and overreliance that can weaken students' own critical thinking \parencite{qian2025}.

These capabilities match the gaps described above. Speech-to-text and summarization can turn a raw lecture recording into searchable, structured material. Text generation can draft individual feedback based on a solution proposal and assessment criteria. Language models can also explain and rank options, such as elective courses, based on a student's stated interests. In all three cases, the system does work that today is either done manually or not done at all. At the same time, the output can sound convincing but still be wrong, and errors in feedback or course advice directly affect students. Human oversight is therefore necessary. This motivates the human-in-the-loop design in Section~\ref{sec:tobe}, the discussion of challenges in Section~\ref{sec:transition} and the sustainability analysis in Section~\ref{sec:susaf}.

In the Gartner Hype Cycle from August 2024 shown in the course lectures, GenAI was at the Peak of Inflated Expectations, with a plateau of productivity forecast in two to five years \parencite{toussaint2025}. Technologies at this stage are often expected to deliver more than they can. This is a reason to be explicit about realistic benefits and limits, and to introduce the solution gradually rather than all at once.

\subsection{An information systems perspective}\label{sec:is}

An information system is a set of interrelated components (hardware, software, databases, networks, people and procedures) that collect, process, store and disseminate information to support decision making, coordination, control, analysis and visualization. It is a sociotechnical system that must fit the organization's culture and goals \parencite{toussaint2025}. Introducing GenAI at NTNU is therefore an organizational change, not just a tool addition. New software alone does not change the process unless the people, procedures and data around it change as well.

Information systems can also be classified by the kind of work they support \parencite{stair}. In the teaching process at NTNU, Studentweb and FS act as transaction processing systems, recording structured, repetitive transactions such as course registrations, exam registrations and grades. Canvas, Panopto and Inspera support the delivery of teaching and assessment, but mainly store and distribute content. None of these systems supports decisions that require interpretation, such as which courses fit a student or what feedback a submission needs. The proposed GenAI service adds this missing layer. A decision support system helps a user make a decision by organizing and analysing relevant information, without making the decision itself. The course recommender in Section~\ref{sec:tobe-case3} works in this way for students and study advisors, while transcription and feedback generation use artificial intelligence to support knowledge work, with humans making the final decisions.

The course distinguishes sustaining (incremental) innovation, which gradually improves existing products, services and ways of operating, from disruptive (radical) innovation, which creates new paradigms. It also distinguishes business process re-engineering (radical, top-down redesign) from continuous improvement (gradual, bottom-up) \parencite{toussaint2025}. The proposed solution is mainly process innovation: the sub-processes are redesigned, modelled in BPMN, and supported by GenAI within existing systems. This is closer to sustaining innovation and continuous improvement than to re-engineering the whole teaching process. It also keeps the risk lower, since the existing systems remain available if the GenAI service fails.

\subsection{Adoption and change}\label{sec:adoption}

Success depends on whether people use the system. The Technology Acceptance Model (TAM) explains use through perceived usefulness and perceived ease of use, and UTAUT adds social influence and facilitating conditions. The DeLone and McLean model links information, system and service quality to use, satisfaction and net benefits. Rogers' diffusion of innovation theory describes uneven adoption across innovators, early adopters, majorities and laggards, and Lewin's change model (unfreezing, moving, refreezing) structures the change process \parencite{stair,toussaint2025}.

These frameworks are used in Section~\ref{sec:transition} to plan the transition and address resistance and literacy risks. TAM explains why the new functionality should be built into tools the staff already use, such as the Canvas grading view, rather than offered as separate systems. Rogers' theory supports starting with pilots among positive lecturers and study programmes before wider rollout. Lewin's model structures each transition from motivating the change to making it the new standard. The DeLone and McLean model provides the basis for evaluating whether the change has succeeded.

\subsection{Sustainability}\label{sec:sustainability}

A new information system affects more than the process it is designed for. The Sustainability Awareness Framework (SusAF) is used to identify such effects across five dimensions (social, individual, environmental, economic and technical) \parencite{duboc2020}, and each effect can be placed on one of three levels: immediate, enabling or structural \parencite{becker2016}. For GenAI, this is relevant because the benefits for students and staff come together with costs that are easy to overlook, such as energy use, processing of personal data and dependency on external vendors. The analysis in Section~\ref{sec:susaf} uses SusAF to assess these effects for the TO-BE solution as a whole and for each case.